\documentclass[11pt]{article}

\usepackage[margin=25mm]{geometry}

\usepackage{amsmath,amssymb,fontspec,unicode-math}

\IfFontExistsTF{DejaVu Serif}{\setmainfont{DejaVu Serif}}{\setmainfont{Latin Modern Roman}}

\IfFontExistsTF{Latin Modern Math}{\setmathfont{Latin Modern Math}}{\setmathfont{DejaVu Math TeX Gyre}}

\usepackage{xurl,hyperref}

\setlength{\parindent}{0pt}\setlength{\parskip}{6pt}

\begin{document}

\title{SWE-SQL: Illuminating LLM Pathways to Solve User SQL Issues in Real-World Applications}

\author{ELVIA preliminary AI analysis}\date{21 October 2026}\maketitle

Preliminary assessment: Minor.\par

Independent scientific review remains pending. Author adoption has not been requested.

Released SQL comparison and normalization discard valid semantic distinctions. A later normalization update fixes some whitespace cases but leaves literal DISTINCT removal; average-tries reporting needs clarification rather than a presumed numeric correction.

\section*{Review scope}

Rapid current-version reassessment of every substantive concern in the supplied LaTeX report: the implicated paper passages and relevant author-released source paths were inspected, existing proposed corrections were briefly checked, and impact was calibrated to the demonstrated scope. No author-provided or historical verification program was executed.

\section*{Problem 1: The result comparator loses semantic distinctions and rejects correct empty results}

Source: Appendices C–D, PDF p. 25; released PostgreSQL utility preprocess\_results, lines 10–29, and ex\_baselines, lines 92–134.

Assessment: supported; impact: minor

\subsection*{Source claim (paraphrase)}

Appendices C–D, PDF p. 25, describe functional SQL evaluation. The released PostgreSQL utility preprocess\_results, lines 10–29, and ex\_base, lines 86–132, discard duplicates/time-of-day and reject empty answers. Required ordering is also not enforced by set comparison.

\subsection*{Analysis}

Set comparison equates the row sequences \(((1),(1))\) and \(((1))\), losing multiplicity and any required order. Rejecting either empty list rejects two identical valid empty answers. Formatting timestamps as dates equates same-day values \(22\) hours apart. These are distinct valid-answer contracts that a general semantic comparator must preserve; their frequency in the final task tests is unknown.

\subsection*{Finding and scope}

This is a supported local finding within the stated source/formulation scope. Its effect on the paper's aggregate experimental results remains unmeasured.

\section*{Problem 2: Removing DISTINCT with text operations changes valid SQL values and query structure}

Source: Released PostgreSQL utility remove\_distinct, lines 32–61; retained conference and current-source versions.

Assessment: supported; impact: minor

\subsection*{Source claim (paraphrase)}

The conference remove\_distinct, lines 32–61 of the PostgreSQL utility, splits each complete SQL string on whitespace, removes tokens equal to distinct, and joins the remainder with single spaces. The current version instead uses a case-insensitive regular expression replacing every whole-word occurrence of DISTINCT with an empty string. The authors’ README explicitly announces an evaluation-code update to this function on \(8\) January 2026.

\subsection*{Analysis}

The current helper applies a case-insensitive whole-word regular expression to the entire SQL string. Thus the valid statement SELECT 'DISTINCT' AS value is changed into SELECT '' AS value. This changes a string value, rather than a SQL keyword. The former split/join behavior that collapsed double spaces or moved a statement into a line comment has been fixed and is excluded from the current defect. Actual benchmark incidence remains unknown.

\subsection*{Finding and scope}

This is a supported local finding within the stated source/formulation scope. Its effect on the paper's aggregate experimental results remains unmeasured.

\section*{Problem 3: The denominator and stopping rule for Table 4’s average tries require clarification}

Source: Section 4, PDF p. 5; Section 6.4, Table 4, PDF p. 9.

Assessment: needs\_evidence; impact: needs\_evidence

\subsection*{Source claim (paraphrase)}

Section 4, PDF p. 5, describes SIX-GYM as containing approximately \(3301\) generated training instances. Section 6.4 and Table 4, p. 9, compare trajectory generation on SIX-GYM. The Reject + f-Plan row has at most five tries, \(2560\) successful trajectories, and an average of \(1.7\) tries. The same table values remain in arXiv v4. The table itself does not explicitly specify its average’s denominator.

Suppose average tries means the mean over all \(3301\) tasks; each successfully solved task consumes at least one try; and each final failure exhausts all five allowed tries. These are an explicit interpretation being tested, not assumptions asserted as supported facts about the run.

\subsection*{Analysis}

Assume the mean covers all \(3301\) tasks, every one of \(2560\) successes takes at least one try, and each of \(741\) failures exhausts five tries. The mean must then be at least \((2560+5\cdot741)/3301\approx1.89791\), inconsistent with \(1.7\). A success-only denominator, early aborts or another averaging unit removes this contradiction. The table does not establish the assumptions, so no incorrect entry or replacement value is asserted.

\subsection*{Finding and scope}

Reporting ambiguity requiring denominator/stopping-rule evidence; no incorrect table entry or replacement number is established.

\subsection*{Proposed correction}

Specify the averaging unit, total attempts, success count and terminal-failure behavior.

A local correction or qualification is proposed where specified; author adoption is not established.

\section*{Limits}

No author-provided program, training procedure, model inference or wet-lab experiment was executed in this review.\par

Historical execution results are retained only with that attribution; a source-level defect alone does not identify the program or configuration used for the published experiments.\par

No author contact, adoption, independent human scientific review or novelty determination is established.\par

The findings apply only to the identified claims, source versions and stated assumptions; they do not certify or invalidate the entire paper.\par

\section*{Source identity}

arXiv:2506.18951v4 (24 January 2026)

\url{https://arxiv.org/pdf/2506.18951v4}

Reviewed PDF SHA-256: \nolinkurl{46d803ed78e06cf51240970f4c7f80dd1ab082adb257afa4fbc0fc54a90dc547}

The exact inspected current paper PDF is identified by its SHA-256. The current publisher PDF was freshly downloaded and is byte-identical to the previously retained proceedings copy. Latest author-manuscript history checked: 2506.18951v4 (last updated 2026-01-24T14:12:49Z). The later dated author revision is used as the primary manuscript; its relevant passages were re-read and locations updated. Current author resource: https://github.com/bird-bench/BIRD-CRITIC-1 at commit f408d6c9db9f16e75b44a1e23f22c04579d1f390.  The selected PDF was freshly downloaded at the recorded retrieval timestamp. Version metadata and operative files were checked separately; file modification times were not used as release dates.

\end{document}
